\documentclass{beamer}
\usepackage[utf8]{inputenc}
\usepackage[T1]{fontenc}
\usepackage{amsmath,amssymb}

\usetheme{Madrid}
\usecolortheme{default}

\title{Your Presentation Title}
\author{Your Name}
\institute{Your Institution}
\date{\today}

\begin{document}

\begin{frame}
  \titlepage
\end{frame}

\begin{frame}{Outline}
  \tableofcontents
\end{frame}

\section{Introduction}

\begin{frame}{Introduction}
  \begin{itemize}
    \item First point about your topic
    \item Second point with more context
    \item Third point leading to your main idea
  \end{itemize}
\end{frame}

\section{Main Content}

\begin{frame}{Key Concepts}
  \begin{block}{Definition}
    A clear definition of your key concept.
  \end{block}

  \begin{alertblock}{Important}
    Something the audience should pay attention to.
  \end{alertblock}

  \begin{exampleblock}{Example}
    A concrete example illustrating the concept.
  \end{exampleblock}
\end{frame}

\begin{frame}{Mathematics}
  You can include equations:
  \begin{equation}
    \int_0^\infty e^{-x^2} \, dx = \frac{\sqrt{\pi}}{2}
  \end{equation}

  Or inline math like $E = mc^2$.
\end{frame}

\section{Conclusion}

\begin{frame}{Conclusion}
  \begin{enumerate}
    \item Summary point one
    \item Summary point two
    \item Summary point three
  \end{enumerate}

  \begin{center}
    \textbf{Thank you! Questions?}
  \end{center}
\end{frame}

\end{document}
